\section{The leWorldLayout Format}

The authored document has one primary object, \texttt{world\_layout}, plus optional environment
metadata. Listing~\ref{lst:minimal-layout} shows the minimal shape.

\begin{figure}[t]
\begin{verbatim}
{
  "schema_version": "le-world-layout-0.1",
  "world_layout": {
    "name": "desk-setup",
    "objects": [],
    "scenario_time_ms": 0,
    "scenario_duration_ms": 0,
    "urdf_xml": "<robot name=\"demo\"/>",
    "joint_positions": {},
    "cameras": []
  },
  "environment": {
    "frame_convention": "ros-rep-103"
  }
}
\end{verbatim}
\caption{Minimal \texttt{leWorldLayout} document shape.}
\label{lst:minimal-layout}
\end{figure}

\subsection{Object Model}

Every object records a stable id, pose, dimensions, type, optional semantic role, optional
appearance representations, optional physics geometry, and optional consistency metadata linking
appearance to physics.

World objects are represented with three layers:
\begin{itemize}[leftmargin=*]
  \item \textbf{identity}: \texttt{id}, \texttt{name}, and optional semantic role;
  \item \textbf{layout}: \texttt{type}, \texttt{position\_xyz}, \texttt{rotation\_rpy\_rad},
    \texttt{size\_xyz}, and \texttt{color};
  \item \textbf{transfer}: \texttt{appearance}, \texttt{physics}, and \texttt{consistency}.
\end{itemize}

\subsection{Appearance and Physics}

The central design choice is that appearance does not imply collision. The \texttt{appearance}
layer stores render and perception representations such as primitives, meshes, and Gaussian splats.
The \texttt{physics} layer stores collision geometry, mass, inertia, friction, restitution, and
fixed/collision flags. The \texttt{consistency} layer records how the two correspond, for example
a hand-authored collider, bounding-box fit, convex decomposition, or unchecked import.

\subsection{Portable Asset References}

Every asset reference must be a portable relative path. References may not include URI schemes,
absolute paths, empty path segments, or \texttt{.}/\texttt{..} traversal. This rule is what lets a
layout travel as a folder or registry artifact without rewriting paths for each machine.

\subsection{Boundary With Scenarios}

A layout is not a task. A scenario references a layout and adds task instruction, role bindings,
randomization, policy and control timing, success conditions, guards, and evaluation metrics. This
separation lets the same world support inspection, editing, simulator transfer, dataset packaging,
and benchmark execution.

